\begin{lemma}
\label{lemma-descend-quasi-finite-universally-open}
Let $S = \lim S_i$ be a limit of a directed system of schemes
with affine transition morphisms.
Let $0 \in I$ and let $f_0 : X_0 \to Y_0$ be a morphism of schemes over $S_0$.
Assume $S_0$, $X_0$, $Y_0$ are quasi-compact and quasi-separated.
Let $f_i : X_i \to Y_i$ be the base change of $f_0$ to $S_i$ and
let $f : X \to Y$ be the base change of $f_0$ to $S$.
If
\begin{enumerate}
\item $f$ is locally quasi-finite and universally open, and
\item $f_0$ is locally of finite presentation,
\end{enumerate}
then there exists an $i \geq 0$ such that $f_i$ is locally quasi-finite
and universally open.
\end{lemma}

\begin{proof}
By Limits, Lemma \ref{limits-lemma-descend-quasi-finite} after increasing
$0$ we may assume $f_0$ is locally quasi-finite. Let $x \in X$.
By \'etale localization of quasi-finite
morphisms we can find a diagram
$$
\xymatrix{
X \ar[d] & U \ar[l] \ar[d] \\
Y & V \ar[l]
}
$$
where $V \to Y$ is \'etale, $U \subset X_V$ is open, $U \to V$ is finite,
and $x$ is in the image of $U \to X$, see
Lemma \ref{lemma-etale-makes-quasi-finite-finite-at-point}.
After shrinking $V$ we may assume $V$ and $U$ are affine.
Since $X$ is quasi-compact, it follows, by taking a finite disjoint
union of such $V$ and $U$, that we can make a diagram as above
such that $U \to X$ is surjective. By
Limits, Lemmas \ref{limits-lemma-descend-finite-presentation},
\ref{limits-lemma-descend-opens},
\ref{limits-lemma-descend-surjective},
\ref{limits-lemma-descend-finite-finite-presentation},
\ref{limits-lemma-descend-etale}, and
\ref{limits-lemma-limit-affine} after possibly increasing $0$
we may assume we have a diagram
$$
\xymatrix{
X_0 \ar[d] & U_0 \ar[l] \ar[d] \\
Y_0 & V_0 \ar[l]
}
$$
where $V_0$ is affine, $V_0 \to Y_0$ is \'etale,
$U_0 \subset (X_0)_{V_0}$ is open, $U_0 \to V_0$ is finite,
and $U_0 \to X_0$ is surjective. Since $V_i \to Y_i$ is \'etale
and hence universally open, follows that it suffices
to prove that $U_i \to V_i$ is universally open for large
enough $i$. This reduces us to the case discussed in the next
paragraph.

\medskip\noindent
Let $A = \colim A_i$ be a filtered colimit of rings.
Let $A_0 \to B_0$ be a ring map. Set $B = A \otimes_{A_0} B_0$ and
$B_i = A_i \otimes_{A_0} B_0$. Assume $A_0 \to B_0$ is finite,
of finite presentation, and $A \to B$ is universally open.
We have to show that $A_i \to B_i$ is universally open for $i$ large enough.
Pick $b_{0, 1}, \ldots, b_{0, d} \in B_0$ which generate $B_0$
as an $A_0$-module. Set $h_0 = \sum_{j = 1, \ldots, d} x_jb_{0, j}$
in $B_0[x_1, \ldots, x_d]$. Denote $h$, resp.\ $h_i$ the image of $h_0$
in $B[x_1, \ldots, x_d]$, resp.\ $B_i[x_1, \ldots, x_d]$.
The image $U$ of $D(h)$ in $\Spec(A[x_1, \ldots, x_d])$ is open
as $A \to B$ is universally open. Of course $U$ is quasi-compact
as the image of an affine scheme. For $i$ large enough there
is a quasi-compact open $U_i \subset \Spec(A_i[x_1, \ldots, x_d])$
whose inverse image in $\Spec(A[x_1, \ldots, x_d])$ is $U$, see
Limits, Lemma \ref{limits-lemma-descend-opens}.
After increasing $i$ we may assume that $D(h_i)$ maps
into $U_i$; this follows from the same lemma by considering the
pullback of $U_i$ in $D(h_i)$. Finally, for $i$ even larger
the morphism of schemes $D(h_i) \to U_i$ will be surjective
by an application of the already used
Limits, Lemma \ref{limits-lemma-descend-surjective}.
We conclude $A_i \to B_i$ is universally open by
Lemma \ref{lemma-characterize-universally-open-finite}.
\end{proof}